% ============================================================
%  Section: Conditional MeshGraphNet for Thermoelastic Wave Prediction
%  Paste this block into your Overleaf document.
%  Requires: amsmath, amssymb, bbm (or dsfont) for \mathbb.
% ============================================================

\section{Conditional MeshGraphNet for Thermoelastic Wave Prediction}

% --------------------------------------------------------
\subsection{Role of MeshGraphNet in the AI Pipeline}
% --------------------------------------------------------

\texttt{ConditionalMeshGraphNet} serves as the graph-based supervised surrogate
model in the AI pipeline.
Its purpose is to learn a fast approximation of the state-transition operator
produced by COMSOL Multiphysics finite-element simulations of coupled
thermoelastic wave propagation in geological media.
The model is not physics-informed in the sense of Ritz-Galerkin or
collocation-based PDE residual minimisation; it is trained end-to-end by
comparing its predictions against COMSOL-generated ground-truth fields.

The learned operator is written compactly as
\[
  \hat{\mathbf{Y}}^{t+\Delta t}
  = f_{\theta}\!\left(\mathbf{X}^{t},\,\mathcal{G}\right),
\]
where \(\mathbf{X}^{t}\) is the matrix of node features at time \(t\),
\(\mathcal{G}\) is the graph representation of the COMSOL mesh,
and \(\hat{\mathbf{Y}}^{t+\Delta t}\in\mathbb{R}^{N\times C}\) collects
the predicted physical fields at the next time step for all \(N\) nodes and
\(C\) output channels.
The physical connection is established implicitly through the mesh topology,
node and edge features, and local message passing, rather than through an
explicit PDE constraint in the loss.

% --------------------------------------------------------
\subsection{Graph Representation of the COMSOL Mesh}
% --------------------------------------------------------

The COMSOL finite-element mesh is represented as a directed graph
\[
  \mathcal{G} = (\mathcal{V},\,\mathcal{E}),
\]
where each node \(v_i\in\mathcal{V}\) corresponds to a mesh node at spatial
position \(\mathbf{p}_i\in\mathbb{R}^{3}\), and each directed edge
\((i,j)\in\mathcal{E}\) represents a local neighbourhood connection.
Graph connectivity is stored in \texttt{edge\_index}:
the first row contains source indices \(i\) and the second row contains
destination indices \(j\), so that information flows from \(v_i\) to \(v_j\).

For each edge the relative displacement vector and its Euclidean length are
\[
  \mathbf{r}_{ij} = \mathbf{p}_j - \mathbf{p}_i,
  \qquad
  d_{ij} = \|\mathbf{r}_{ij}\|_2,
\]
and the edge feature vector is
\[
  \mathbf{a}_{ij}
  =
  \begin{bmatrix}
    \mathbf{r}_{ij} \\ d_{ij}
  \end{bmatrix}.
\]
This encoding preserves the local geometry of the discretisation without
projecting the unstructured mesh onto a regular Cartesian grid.

% --------------------------------------------------------
\subsection{Node and Edge Features}
% --------------------------------------------------------

At each time step \(t\), the feature vector of node \(i\) concatenates four
groups of information:
\[
  \mathbf{x}_i^t
  =
  \begin{bmatrix}
    \mathbf{p}_i \\[2pt]
    \mathbf{m}_i \\[2pt]
    \mathbf{s}   \\[2pt]
    \mathbf{q}_i^t
  \end{bmatrix},
\]
where \(\mathbf{p}_i\) are spatial coordinates, \(\mathbf{m}_i\) are
per-node material parameters (thermal conductivity, density, heat capacity,
Young's modulus, Poisson's ratio, thermal expansion coefficient),
\(\mathbf{s}\) is a scenario-level conditioning vector (source location,
temperature, heating radius, and other global simulation parameters),
and \(\mathbf{q}_i^t\) is the current dynamic state comprising temperature,
displacement, velocity, stress, and strain components.

Critically, this concatenation is performed by the data pipeline before the
model receives any input; the model itself exposes a single unified
\texttt{node\_features} tensor and does not branch on condition separately.
The full node feature matrix is
\(\mathbf{X}^{t}\in\mathbb{R}^{N\times F_{v}}\)
and the edge attribute matrix is
\(\mathbf{A}\in\mathbb{R}^{M\times F_{e}}\),
where \(M=|\mathcal{E}|\) is the number of directed edges.

% --------------------------------------------------------
\subsection{Prediction Target and Delta Formulation}
% --------------------------------------------------------

The model supports two target modes, controlled by \texttt{target\_mode} in
the dataset metadata.

\paragraph{Absolute mode.}
The network directly predicts the next full state:
\(\hat{\mathbf{Y}}^{t+\Delta t} = f_{\theta}(\mathbf{X}^{t},\mathcal{G})\).

\paragraph{Delta mode (default).}
The network predicts the state increment for each node:
\[
  \Delta\hat{\mathbf{q}}_i^t
  = f_{\theta}\!\left(\mathbf{x}_i^t,\,\mathcal{G}\right),
\]
and the next state is reconstructed by
\[
  \hat{\mathbf{q}}_i^{t+\Delta t}
  = \mathbf{q}_i^t + \Delta\hat{\mathbf{q}}_i^t.
\]
Delta prediction is generally preferable for thermoelastic trajectories because
inter-frame increments are smoother than the absolute fields, which facilitates
learning and reduces the magnitude of supervised targets.

During autoregressive inference the predicted state replaces the dynamic part
of the node features and is fed back as input for the next step:
\[
  \hat{\mathbf{q}}^{t+(r+1)\Delta t}
  = \hat{\mathbf{q}}^{t+r\Delta t}
    + f_{\theta}\!\left(\hat{\mathbf{X}}^{t+r\Delta t},\,\mathcal{G}\right),
  \qquad r = 0,1,\dots,R-1,
\]
where \(\hat{\mathbf{X}}^{t+r\Delta t}\) contains static features
(\(\mathbf{p}_i,\mathbf{m}_i,\mathbf{s}\)) together with the
predicted dynamic state from the previous rollout step.

% --------------------------------------------------------
\subsection{Encoder--Processor--Decoder Architecture}
% --------------------------------------------------------

The overall model follows the Encoder--Processor--Decoder decomposition:
\[
  f_{\theta}
  = D_{\theta_D} \circ P_{\theta_P} \circ E_{\theta_E}.
\]

\paragraph{Encoder.}
Separate MLPs project the raw node features and edge attributes into a shared
latent space of dimension \(H\):
\[
  \mathbf{h}_i^0 = \phi_v^{\mathrm{enc}}(\mathbf{x}_i^t),
  \qquad
  \mathbf{e}_{ij}^0 = \phi_e^{\mathrm{enc}}(\mathbf{a}_{ij}).
\]

\paragraph{Processor.}
A \texttt{ModuleList} of \(K\) \texttt{MessagePassingBlock} modules
iteratively refines the latent node and edge representations (detailed in
Section~\ref{sec:message_passing}).

\paragraph{Decoder.}
A final MLP without output layer normalisation maps each node's latent vector
to the predicted output channels:
\[
  \hat{\mathbf{y}}_i = \phi^{\mathrm{dec}}(\mathbf{h}_i^K),
\]
yielding the output matrix
\(\hat{\mathbf{Y}} = [\hat{\mathbf{y}}_1,\dots,\hat{\mathbf{y}}_N]^{\top}
\in\mathbb{R}^{N\times C}\).

\paragraph{MLP design.}
Every MLP in the encoder, processor, and decoder is built by
\texttt{make\_mlp}: a sequence of \texttt{Linear} layers with
optional \texttt{LayerNorm}, \texttt{SiLU} activation, and \texttt{Dropout}
in all hidden layers.
The decoder MLP omits the output \texttt{LayerNorm} to allow unconstrained
regression.
A single hidden-layer transformation reads
\[
  \mathbf{z}_{\ell+1}
  = \operatorname{Dropout}\!\Bigl(
      \operatorname{SiLU}\!\bigl(
        \operatorname{LayerNorm}(\mathbf{W}_\ell\mathbf{z}_\ell + \mathbf{b}_\ell)
      \bigr)
    \Bigr),
  \qquad
  \operatorname{SiLU}(x) = x\cdot\sigma(x).
\]
Default hyperparameters: \(H=128\), \(K=10\), \texttt{mlp\_layers}=3,
\texttt{dropout}=0.05, \texttt{layer\_norm}=\texttt{True}.

% --------------------------------------------------------
\subsection{Message Passing Mechanism}
\label{sec:message_passing}
% --------------------------------------------------------

Each \texttt{MessagePassingBlock} processes the latent state in three stages.

\paragraph{Edge update.}
For each edge \((i,j)\) the updated edge latent vector is computed from the
concatenation of the source node vector, the destination node vector, and the
current edge vector:
\[
  \tilde{\mathbf{e}}_{ij}^{k+1}
  = \phi_e^k\!\left(
      \bigl[\mathbf{h}_i^k,\;\mathbf{h}_j^k,\;\mathbf{e}_{ij}^k\bigr]
    \right),
\]
followed by a residual connection and layer normalisation:
\[
  \mathbf{e}_{ij}^{k+1}
  = \operatorname{LayerNorm}\!\left(
      \mathbf{e}_{ij}^{k} + \tilde{\mathbf{e}}_{ij}^{k+1}
    \right).
\]
The input to \(\phi_e^k\) has dimension \(3H\) (\texttt{edge\_mlp}); the
output has dimension \(H\).

\paragraph{Message aggregation.}
Incoming messages for each destination node are summed via a
\texttt{scatter\_add} over the destination indices:
\[
  \mathbf{m}_j^{k+1}
  = \sum_{i:\,(i,j)\in\mathcal{E}} \mathbf{e}_{ij}^{k+1}.
\]

\paragraph{Node update.}
The aggregated message is concatenated with the current node latent vector,
processed by a second MLP, and stabilised with a residual connection:
\[
  \tilde{\mathbf{h}}_j^{k+1}
  = \phi_v^k\!\left(
      \bigl[\mathbf{h}_j^k,\;\mathbf{m}_j^{k+1}\bigr]
    \right),
  \qquad
  \mathbf{h}_j^{k+1}
  = \operatorname{LayerNorm}\!\left(
      \mathbf{h}_j^k + \tilde{\mathbf{h}}_j^{k+1}
    \right).
\]
The input to \(\phi_v^k\) has dimension \(2H\) (\texttt{node\_mlp}).

After \(K\) steps, the latent vector \(\mathbf{h}_i^K\) aggregates
information from the \(K\)-hop neighbourhood
\(\mathcal{N}_K(i)=\{j\mid\operatorname{dist}_{\mathcal{G}}(i,j)\le K\}\).
This is physically meaningful because thermoelastic wave propagation is governed
by local interactions: temperature gradients, elastic coupling, and stress
redistribution all depend on the state of neighbouring nodes.

A technical remark: under CUDA automatic mixed precision (AMP), the
scatter buffer is allocated in the dtype of the edge tensor (\texttt{float16}),
and is cast back to the node dtype (\texttt{float32}) before concatenation,
ensuring numerical stability of the residual additions.

% --------------------------------------------------------
\subsection{Loss Function and Training Procedure}
% --------------------------------------------------------

\paragraph{WeightedFieldMSE.}
The model is trained with a weighted mean squared error that assigns
independent weights to physical field groups:
\[
  \mathcal{L}
  = \frac{1}{\sum_{c=1}^{C} w_c}
    \sum_{c=1}^{C} w_c
    \cdot \frac{1}{N}
    \sum_{i=1}^{N}
    \left(\hat{y}_{i,c} - y_{i,c}\right)^2,
\]
where \(\hat{y}_{i,c}\) is the model prediction and \(y_{i,c}\) is the
COMSOL target for field \(c\) at node \(i\).
The weight \(w_c = w_g\) is selected according to the physical group
\(g\in\{T,\,u,\,\dot{u},\,\sigma,\,\varepsilon,\,\mathrm{other}\}\)
to which field \(c\) belongs:
\[
  \mathcal{C}
  = \mathcal{C}_{T}
    \cup \mathcal{C}_{u}
    \cup \mathcal{C}_{\dot{u}}
    \cup \mathcal{C}_{\sigma}
    \cup \mathcal{C}_{\varepsilon}
    \cup \mathcal{C}_{\mathrm{other}}.
\]
Group assignment follows the naming convention of exported COMSOL fields.
The loss contains no PDE residual term; the model learns to approximate
COMSOL transitions in a purely supervised manner.

\paragraph{Optimiser and regularisation.}
Training uses the Adam optimiser over the trainable parameter subset
\(\{p\in\theta \mid p\texttt{.requires\_grad}\}\).
At every gradient step the global gradient norm is clipped:
\[
  \mathbf{g}
  \leftarrow
  \mathbf{g}\cdot\min\!\left(1,\,\frac{\tau}{\|\mathbf{g}\|_2}\right),
\]
where \(\tau\) is the clipping threshold (\texttt{grad\_clip}).
When a CUDA device is available, automatic mixed precision (AMP) is enabled
via \texttt{torch.amp.autocast} and \texttt{GradScaler}.

\paragraph{Validation and early stopping.}
Training and validation losses are tracked epoch-by-epoch.
The checkpoint saved at the epoch with the lowest validation loss is used for
all subsequent evaluation and inference:
\[
  e^{*} = \arg\min_{e}\;\mathcal{L}_{\mathrm{val}}^{(e)}.
\]
Early stopping halts training when the validation loss does not improve for a
configurable patience window.
Each checkpoint stores the model state dict, optimiser state, epoch index,
validation loss, field names, input/output dimensions, dataset identifiers,
model configuration, and target mode.

% --------------------------------------------------------
\subsection{Fine-Tuning Strategy}
% --------------------------------------------------------

Three fine-tuning modes are supported via \texttt{freeze\_for\_finetune}:

\begin{itemize}
  \item \textbf{full}: all parameters remain trainable,
        \(\theta_{\mathrm{train}} = \{\theta_E,\theta_P,\theta_D\}\).
        Used when the new dataset differs substantially from the pretraining
        distribution.

  \item \textbf{decoder\_only}: encoder and processor are frozen,
        \(\theta_{\mathrm{train}} = \{\theta_D\}\).
        Effective when the latent dynamics are transferable but the output
        mapping must be recalibrated (e.g.\ different field normalisation or
        new rock type with the same feature schema).

  \item \textbf{processor\_decoder}: only the encoder is frozen,
        \(\theta_{\mathrm{train}} = \{\theta_P,\theta_D\}\).
        Used when the input-feature structure is consistent across datasets
        but local interaction patterns change (e.g.\ different anisotropy or
        boundary configuration).
\end{itemize}

Fine-tuning is invoked through \texttt{fine\_tune\_model}, which delegates to
\texttt{train\_model} with the \texttt{fine\_tune=True} flag and reads the
mode from \texttt{config[``fine\_tune''][``mode'']}.

% --------------------------------------------------------
\subsection{Validation and Physical Evaluation}
% --------------------------------------------------------

The validation pipeline in \texttt{run\_full\_validation} operates at four
levels.

\paragraph{One-step validation.}
For every sample \((\mathbf{X}^t, \mathbf{Y}^{t+\Delta t})\) in the
evaluation split, the model predicts the next state from the true current
state.
Per-field RMSE and MAE are reported:
\[
  \mathrm{RMSE}_c
  = \sqrt{\frac{1}{N}\sum_{i=1}^{N}\!\left(\hat{y}_{i,c}-y_{i,c}\right)^2},
  \qquad
  \mathrm{MAE}_c
  = \frac{1}{N}\sum_{i=1}^{N}\left|\hat{y}_{i,c}-y_{i,c}\right|.
\]

\paragraph{Autoregressive rollout validation.}
Starting from the first sample of the split, the model is applied
recurrently, feeding each predicted state back as input.
Rollout RMSE over time is tracked per field, making error accumulation
visible.

\paragraph{Derived physical metrics.}
The following quantities are computed from the raw trajectories for both the
COMSOL reference and the model prediction:

\begin{itemize}
  \item displacement magnitude:
        \(\|\mathbf{u}_i\|_2 = \sqrt{u_i^2+v_i^2+w_i^2}\);
  \item velocity magnitude:
        \(\|\dot{\mathbf{u}}_i\|_2 = \sqrt{\dot{u}_i^2+\dot{v}_i^2+\dot{w}_i^2}\);
  \item temperature change:
        \(\Delta T_i^t = T_i^t - T_i^{0}\);
  \item von Mises stress (from the exported \texttt{solid.mises} field if
        available, otherwise reconstructed as
        \[
          \sigma_{\mathrm{vM}}
          = \sqrt{
              \tfrac{1}{2}
              \bigl[
                (\sigma_x\!-\!\sigma_y)^2
                +(\sigma_y\!-\!\sigma_z)^2
                +(\sigma_z\!-\!\sigma_x)^2
                +6(\sigma_{xy}^2+\sigma_{xz}^2+\sigma_{yz}^2)
              \bigr]
            }
        \,\text{);}
  \]
  \item stress risk flag:
        \(R_i^t = \mathbb{I}\!\left(
            \sigma_{\mathrm{vM},i}^t / \max_j \sigma_{\mathrm{vM},j}^t
            \ge 0.8
          \right)\).
\end{itemize}

\paragraph{Wave-front radius.}
The thermal wave front is estimated from the thresholded temperature-change
field.
Let \(\mathbf{p}_{\mathrm{src}}\) be the source centre derived from
scenario metadata and let \(\mathcal{A}^t\) be the set of active nodes.
The front radius is approximated as the 95th percentile of distances from
active nodes to the source:
\[
  R_{\mathrm{front}}^t
  = \operatorname{Percentile}_{95}\!\left(
      \left\{
        \|\mathbf{p}_i - \mathbf{p}_{\mathrm{src}}\|_2
        \mid i\in\mathcal{A}^t
      \right\}
    \right).
\]
Both COMSOL and predicted front radii are tabulated over time, and their
mean absolute error is recorded in the validation summary.

\paragraph{Visual reports.}
For a selected set of time steps the validation pipeline generates comparison
slices showing the COMSOL reference field, the model prediction, and the
pointwise absolute error \(|\hat{y}_{i,c}^t - y_{i,c}^t|\).
Additionally, bar charts of the worst fields by one-step and final-rollout
relative RMSE, line plots of RMSE over rollout steps for raw and derived
fields, and an HTML summary report are written to the output directory.

% --------------------------------------------------------
\subsection{Advantages and Limitations}
% --------------------------------------------------------

\paragraph{Advantages.}
The primary advantage of \texttt{ConditionalMeshGraphNet} is its native
support for unstructured COMSOL meshes: unlike grid-based methods, no
interpolation to a regular lattice is required.
Material and scenario conditioning is incorporated without architectural
modifications, because all conditioning information is concatenated into the
node feature vector by the data pipeline.
Local message passing provides an inductive bias aligned with
the locality of thermoelastic interactions.
After training, inference requires only a single forward pass per time step,
making the model substantially faster than a full COMSOL solve for
repeated or real-time prediction scenarios.

\paragraph{Limitations.}
The model is bounded by the COMSOL training distribution.
For inputs far outside the training set,
\((\mathbf{X}_{\mathrm{new}},\mathcal{G}_{\mathrm{new}})\not\sim\mathcal{D}_{\mathrm{train}}\),
predictions may deteriorate without warning.
Conservation laws, constitutive relations, and boundary conditions are learned
implicitly from data and are not enforced algebraically during inference.
Autoregressive rollout accumulates prediction errors, so long-horizon
predictions must be validated independently of one-step metrics.
Training memory scales as
\(\mathcal{M} \propto (N+M)H + KMH\),
where \(K\) is the number of message passing blocks and \(H\) the latent
dimension, which may limit the mesh resolution for large 3D domains.
Finally, strict consistency between training and inference is required:
field ordering, normalisation statistics, graph construction, edge attributes,
and target mode must be identical at every stage of the pipeline.

% --------------------------------------------------------
\subsection{Summary}
% --------------------------------------------------------

\texttt{ConditionalMeshGraphNet} was implemented as a supervised graph-based
surrogate for thermoelastic wave prediction on COMSOL finite-element meshes.
The mesh is represented as \(\mathcal{G}=(\mathcal{V},\mathcal{E})\), where
nodes store geometry, material parameters, scenario conditioning, and the
current physical state, while edges store relative geometric information.
The node feature vector
\[
  \mathbf{x}_i^t
  =
  \begin{bmatrix}
    \mathbf{p}_i \\ \mathbf{m}_i \\ \mathbf{s} \\ \mathbf{q}_i^t
  \end{bmatrix}
\]
is assembled by the data pipeline; no separate conditioning branch is needed.
The model learns the COMSOL state-transition approximation
\(\hat{\mathbf{Y}}^{t+\Delta t}=f_{\theta}(\mathbf{X}^{t},\mathcal{G})\)
via the composition
\(f_{\theta} = D_{\theta_D}\circ P_{\theta_P}\circ E_{\theta_E}\).
The processor applies \(K=10\) \texttt{MessagePassingBlock} layers, each
performing edge update, \texttt{scatter\_add} aggregation, and node update
with residual connections.
Training minimises \texttt{WeightedFieldMSE} over COMSOL targets, grouped
into temperature, displacement, velocity, stress, strain, and other fields,
using Adam with gradient clipping, AMP on CUDA, early stopping, and
best-checkpoint selection.

The model should therefore be interpreted as a fast learned approximation of
COMSOL state transitions.
Physical relevance is established through mesh topology, material and scenario
conditioning, and localised message passing; the governing thermoelastic PDE
is not explicitly solved or residualised during training.

% --------------------------------------------------------
\paragraph*{Recommended figures for this section.}
% --------------------------------------------------------
\begin{enumerate}
  \item \textbf{COMSOL mesh to graph representation.}
        Side-by-side view of finite-element nodes, graph nodes
        \(\mathcal{V}\), and graph edges \(\mathcal{E}\) with relative edge
        vectors \(\mathbf{r}_{ij}\).
  \item \textbf{Encoder--Processor--Decoder architecture.}
        Block diagram showing the node encoder, edge encoder, \(K\)
        \texttt{MessagePassingBlock} modules, and the decoder, with tensor
        dimensions annotated.
  \item \textbf{MessagePassingBlock internals.}
        Data-flow diagram of the edge MLP update, \texttt{scatter\_add}
        aggregation, and node MLP update with residual connections and
        \texttt{LayerNorm}.
  \item \textbf{Training and validation pipeline.}
        Flowchart from COMSOL trajectory export to graph dataset construction,
        supervised training, one-step validation, and autoregressive rollout
        evaluation.
  \item \textbf{COMSOL vs.\ prediction vs.\ absolute error slice.}
        Three-panel figure (COMSOL reference \(\mid\) model prediction \(\mid\)
        pointwise absolute error \(|\hat{y}_{i,c}^t-y_{i,c}^t|\)) for
        temperature change, displacement magnitude, or von Mises stress
        at a representative rollout time step.
\end{enumerate}